
\documentclass[12pt, a4paper]{article}

% --- Packages ---
\usepackage[margin=1in]{geometry}
\usepackage{amsmath, amssymb}
\usepackage{hyperref}
\usepackage{xcolor}
\usepackage{titlesec}
\usepackage{enumitem}
\usepackage{tcolorbox}
\usepackage{booktabs}
\usepackage{tabularx}

% --- Formatting ---
\hypersetup{
    colorlinks=true,
    linkcolor=blue,
    filecolor=magenta,      
    urlcolor=teal,
    pdftitle={Literature Review: ZKP-Enabled eKYC Systems},
}

\titleformat{\section}{\large\bfseries\color{darkgray}}{\thesection}{1em}{}[\titlerule]
\titleformat{\subsection}{\normalsize\bfseries\color{gray}}{\thesubsection}{1em}{}

% --- Document Metadata ---
\title{\textbf{Literature Review: Privacy-Preserving \\ e-KYC and Decentralized Identity}}
\author{
    \textbf{Practice School-I (PS-1) Technical Log - Day 1} \\ 
    CDAC Pune \\
    \textit{Topic: Core Zero-Knowledge Cryptography \& Identity Verification}
}
\date{May 29, 2026}
\begin{document}

\maketitle

\begin{abstract}
This document synthesizes four foundational academic papers regarding the application of Zero-Knowledge Proofs (ZKPs) to electronic Know Your Customer (e-KYC) and decentralized identity (DID) frameworks. The review categorizes the current research into selective disclosure mechanics, non-transferable reputation systems, stateless authentication, and off-chain zkVM attestations. This literature review serves as the architectural baseline for the development of a privacy-preserving Indian e-KYC prototype.
\end{abstract}

\vspace{0.5cm}

% ==========================================
% PAPER 1: Selective Disclosure & Merkle Trees
% ==========================================
\section{Selective Disclosure via Merkle Trees and zk-SNARKs}
\textbf{Core Paper:} \textit{A Privacy-Preserving Selectively Disclosed eKYC System Using Merkle Tree and Zero-Knowledge Proofs} (Ahmed et al.)

Traditional centralized databases present severe vulnerabilities to data breaches and unauthorized surveillance. To counter this, Ahmed et al. propose a system that enforces data minimization through "selective disclosure".
\begin{itemize}
    \item \textbf{Merkle Tree Integration:} Instead of storing raw identity data, the system structures user attributes (e.g., Name, Date of Birth) as leaf nodes in a Merkle Tree, obfuscated using a Poseidon hash function alongside cryptographic salts. 
    \item \textbf{Single Proof Generation:} The issuer digitally signs the Merkle Root and commits it to the blockchain. The user can then generate a single Groth16 zk-SNARK proof demonstrating possession of specific attributes (like Age $\ge$ 18) without revealing the entire credential document.
    \item \textbf{Performance Context:} The paper benchmarks proof generation at approximately 312 milliseconds and on-chain verification at 20,000 gas, proving its viability for scalable financial applications.
\end{itemize}

% ==========================================
% PAPER 2: P2P Marketplaces & Soulbound Tokens
% ==========================================
\section{Decentralized Reputation and Soulbound Tokens (SBTs)}
\textbf{Core Paper:} \textit{ZKP Enabled Identity and Reputation Verification in P2P Marketplaces} (Kalbantner et al.)

In decentralized and Peer-to-Peer (P2P) financial marketplaces, maintaining regulatory compliance (AML/CTF) while preserving transactional privacy is a primary challenge.
\begin{itemize}
    \item \textbf{Identity Anchoring:} The authors integrate Verifiable Credentials (VCs) with Soulbound Tokens (SBTs). SBTs are non-transferable NFTs that securely bind a user's verified KYC status directly to their Decentralized Identifier (DID).
    \item \textbf{Tiered Reputation System:} To counter Sybil attacks and incentivize good behavior, the framework introduces Marketplace Reputation Tokens (MRTs). These tokens establish a 3-tier reputation scoring mechanism based on transaction history and varying depths of KYC verification.
    \item \textbf{State Channels:} Trades are initialized through off-chain state channels where ZKPs are used to verify the SBTs and MRTs without publicly exposing the users' actual identities on the ledger.
\end{itemize}

% ==========================================
% PAPER 3: Stateless Authentication & Recovery
% ==========================================
\section{Stateless Authentication via Human-Readable UIDs}
\textbf{Core Paper:} \textit{ZKPass - A Zero-Knowledge Proof Based Authentication System for Secure Identity Verification} (Srinivas P M et al.)

A critical barrier to mainstream Web3 adoption is the complexity of cryptographic key management and the vulnerability of password-based logins. ZKPass proposes a lightweight, passwordless alternative.
\begin{itemize}
    \item \textbf{Ephemeral Key Generation:} The framework derives temporary private keys locally from a Unique Identifier (UID) exclusively at the moment of login. This eliminates the need to store long-term keys or passwords, drastically reducing the attack surface.
    \item \textbf{Self-Sovereign Recovery:} Moving away from centralized recovery agents, the system generates a human-readable English recovery phrase from the UID and a secret salt during registration.
    \item \textbf{Storage Efficiency:} By storing only the hash of the UID on the blockchain rather than complete credential proofs or public keys, the system mitigates blockchain congestion and ensures strict identity privacy.
\end{itemize}

% ==========================================
% PAPER 4: Off-Chain Attestations & zkVM
% ==========================================
\section{Off-Chain Attestations via EAS and zkVM}
\textbf{Core Paper:} \textit{Zero-Knowledge Identity Verification} (Kumar et al.)

Publishing attestations directly on-chain inherently exposes relational metadata, violating privacy-by-default principles. Kumar et al. decouple the attestation issuance from on-chain verification.
\begin{itemize}
    \item \textbf{EIP-712 Standard:} The system relies on the Ethereum Attestation Service (EAS) to issue off-chain, structured attestations signed securely under the EIP-712 standard.
    \item \textbf{zkVM Processing:} Instead of writing bespoke arithmetic circuits, the attestations and their logical constraints are processed inside a Zero-Knowledge Virtual Machine (specifically SP1).
    \item \textbf{Client-Side Verification:} The zkVM compiles succinct Groth16 or Plonk proofs that are lightweight enough to be verified completely client-side using WebAssembly (WASM) in a standard web browser or Node.js environment.
\end{itemize}

% ==========================================
% SYNTHESIS & ARCHITECTURAL ROADMAP
% ==========================================
\section{Synthesis for PS-1 Project Implementation}
The literature confirms that standard on-chain KYC systems are structurally inadequate due to data exposure and gas inefficiencies. To build the required India-centric e-KYC prototype for this Practice School project, the architecture will synthesize three core mechanisms from the reviewed literature:
\begin{enumerate}
    \item \textbf{The Data Structure:} We will adopt the \textit{Merkle Tree} approach (Ahmed et al.) to hash the individual fields of a DigiLocker JSON credential, allowing for selective disclosure (e.g., Age).
    \item \textbf{The Proof Generation:} To ensure broad compatibility and reduced manual circuit errors, we will explore generating the proofs via a \textit{zkVM environment} (Kumar et al.) rather than raw Circom compilation.
    \item \textbf{The Storage Layer:} To minimize blockchain footprint, we will adopt the \textit{Stateless UID protocol} (Srinivas et al.), ensuring the smart contract verifies the zk-SNARK proof without storing any persistent state beyond a single cryptographic hash.
\end{enumerate}

\end{document}
